%% Worksheets, AP Precalculus Unit 1. Block ids end with the worksheet letter (1.2A) or MEGA.
\newcommand{\gr}[2]{\begin{center}\begin{tikzpicture}\begin{axis}[gridaxis,#1]#2\end{axis}\end{tikzpicture}\end{center}}
\newcommand{\pt}[1]{\node[circle,fill=apblue,inner sep=1.5pt] at (axis cs:#1){};}
\newcommand{\op}[1]{\node[circle,draw=apblue,fill=white,thick,inner sep=1.5pt] at (axis cs:#1){};}
\newcommand{\va}[1]{\draw[apred,dashed,thick](axis cs:#1,-100)--(axis cs:#1,100);}
\newcommand{\ha}[1]{\draw[apgold!80!black,dashed,thick](axis cs:-100,#1)--(axis cs:100,#1);}

%%DOC SC1.1-1.3B | Quiz B \textbullet\ Topics 1.1 to 1.3: Rates of Change in Functions
%%PRE \let\apolddirections\directions\let\endapolddirections\enddirections
%%PRE \RenewEnviron{directions}{\begin{apolddirections}Answer every question. Multiple-choice questions have one best answer. On free-response questions, show your work and justify your answers. \textbf{Graphing calculators are not allowed.} A simple (non-graphing) calculator may be used if necessary. For questions with a graph, read the values you need from the graph.\end{apolddirections}}
%%PRE \newcommand{\pts}[1]{{\sffamily\small\bfseries\color{apnavy}#1 pts}}
%%PRE \let\apoldband\titleband
%%PRE \renewcommand{\titleband}[2]{{\centering\sffamily\bfseries\Large\color{apnavy}The Senegalese American Bilingual School\par}\vspace{6pt}\apoldband{#1}{Class: 11th B \quad #2\ifkey\else\quad Score: \rule{1.4cm}{0.4pt}\,/\,100\fi\par\vspace{2pt}{\small Duration: 120 minutes \textbullet\ Simple calculator only, no graphing calculator \textbullet\ 100 points}}}
%%PRE \geometry{a4paper,margin=1.5cm,top=1.6cm,bottom=1.6cm,headheight=14pt}
%%PRE \setlength{\columnsep}{0.8cm}\setlength{\columnseprule}{0.3pt}
\renewcommand{\nocalc}{}
\renewenvironment{exo}[2][]{\refstepcounter{exo}\par\Needspace{5\baselineskip}\vspace{3pt}\noindent{\sffamily\bfseries\color{apnavy}Question \theexo}\hfill\ifx&#2&\else#2\fi\par\nopagebreak}{\par}
\renewcommand{\choices}[4]{\par\vspace{1pt}\begin{tabular}{@{}l@{\hspace{4pt}}p{0.89\linewidth}@{}}(A)&#1\\(B)&#2\\(C)&#3\\(D)&#4\end{tabular}\par}
\renewcommand{\gr}[2]{\begin{center}\begin{tikzpicture}\begin{axis}[gridaxis,#1]#2\end{axis}\end{tikzpicture}\end{center}\vspace{-6pt}}
\newcommand{\side}[2]{\par\noindent\begin{minipage}[c]{0.47\linewidth}#1\end{minipage}\hfill\begin{minipage}[c]{0.5\linewidth}#2\end{minipage}\par\medskip}
\providecommand{\sk}[2]{\begin{center}\begin{tikzpicture}\begin{axis}[axis lines=middle,xtick=\empty,ytick=\empty,xlabel={$x$},ylabel={$y$},width=5.4cm,height=4.4cm,enlargelimits=false,clip=false,#1]#2\end{axis}\end{tikzpicture}\end{center}}
\providecommand{\lab}[4]{\node[circle,fill=black,inner sep=1.4pt] at (axis cs:#1,#2){};\node[font=\small\bfseries,anchor=#4,inner sep=2.5pt] at (axis cs:#1,#2){#3};}
\begin{center}\textsf{\bfseries\color{apnavy}Part A \textbullet\ Free response \textbullet\ No graphing calculator \textbullet\ 52 points}\end{center}

\Needspace{10\baselineskip}
\side{\gr{xmin=-6.5,xmax=6.5,ymin=-4,ymax=4,width=7.4cm,height=4.3cm}{\addplot[apblue,very thick,domain=-6:6,samples=150]{3*cos(deg(pi*x/3))};\pt{-6,3}\pt{6,3}\node[font=\small,apblue,fill=white,inner sep=1pt] at (axis cs:-1.8,-3.4){$y=f(x)$};}}{\begin{exo}[Free response]{\pts{8}}
The graph of a function $f$ is shown for $-6\le x\le6$.
\begin{enumerate}[label=(\alph*)]
\item On what open intervals is $f$ increasing?
\item On what open intervals is $f$ decreasing?
\item On what open intervals is $f$ positive?
\item On what open intervals is $f$ negative?
\end{enumerate}
\ifkey\else\vspace{1.2cm}\fi
\begin{solution}
\textit{Scoring: 2 points per part.}\\
(a) $(-3,0)$ and $(3,6)$. \quad (b) $(-6,-3)$ and $(0,3)$. \quad (c) $(-6,-4.5)$, $(-1.5,1.5)$ and $(4.5,6)$. \quad (d) $(-4.5,-1.5)$ and $(1.5,4.5)$.
\end{solution}
\end{exo}}

\Needspace{10\baselineskip}
\side{\gr{xmin=-0.5,xmax=8.5,ymin=-0.5,ymax=7,width=7.2cm,height=4.3cm}{\addplot[apblue,very thick] coordinates {(0,5) (2,1) (4,1)};\addplot[apblue,very thick,domain=4:7,samples=60]{1+(5/9)*(x-4)^2};\addplot[apblue,very thick] coordinates {(7,6) (8,4)};\pt{0,5}\pt{8,4}\node[font=\small,apblue,fill=white,inner sep=1pt] at (axis cs:1.6,6.2){$y=g(x)$};}}{\begin{exo}[Free response]{\pts{6}}
The graph of a function $g$ is shown for $0\le x\le8$. It is made of line segments and a parabola.
\begin{enumerate}[label=(\alph*)]
\item On what open intervals is $g$ increasing?
\item On what open intervals is $g$ decreasing?
\item What is the rate of change of $g$ on the interval $(2,4)$? Explain.
\end{enumerate}
\ifkey\else\vspace{1cm}\fi
\begin{solution}
\textit{Scoring: 2 points per part.}\\
(a) $(4,7)$. \quad (b) $(0,2)$ and $(7,8)$. \quad (c) $0$: $g$ is constant on $(2,4)$, so it is neither increasing nor decreasing there.
\end{solution}
\end{exo}}

\Needspace{10\baselineskip}
\side{\gr{xmin=-6.5,xmax=6.5,ymin=-4,ymax=4,width=7.4cm,height=4.3cm}{\addplot[apblue,very thick,domain=-6:6,samples=150]{3*sin(deg(pi*x/3))};\node[font=\small,apblue,fill=white,inner sep=1pt] at (axis cs:-1.5,3.3){$y=h(x)$};}}{\begin{exo}[Free response]{\pts{4}}
The graph of a function $h$ is shown for $-6\le x\le6$.
\begin{enumerate}[label=(\alph*)]
\item On what open intervals is the graph of $h$ concave up?
\item On what open intervals is the graph of $h$ concave down?
\end{enumerate}
\ifkey\else\vspace{1cm}\fi
\begin{solution}
\textit{Scoring: 2 points per part.}\\
(a) $(-3,0)$ and $(3,6)$. \quad (b) $(-6,-3)$ and $(0,3)$. The concavity changes at $x=-3$, $0$ and $3$.
\end{solution}
\end{exo}}

\Needspace{14\baselineskip}


\Needspace{12\baselineskip}
\begin{exo}[Free response]{\pts{8}}
For each situation, decide whether the second quantity has a \textbf{positive} or a \textbf{negative} rate of change with respect to the first.
\begin{enumerate}[label=(\alph*),itemsep=2pt]
\item The number of hours a student studies, and the number of mistakes on the test. \ifkey\else\hfill\rule{2.2cm}{0.4pt}\fi
\item The number of friends sharing a taxi fare, and the amount each one pays. \ifkey\else\hfill\rule{2.2cm}{0.4pt}\fi
\item The time since a pot of water was put on the fire, and the temperature of the water. \ifkey\else\hfill\rule{2.2cm}{0.4pt}\fi
\item The number of pages printed, and the ink left in the cartridge. \ifkey\else\hfill\rule{2.2cm}{0.4pt}\fi
\item The radius of a pizza, and its area. \ifkey\else\hfill\rule{2.2cm}{0.4pt}\fi
\item The speed of a bus, and the time it takes to go from Dakar to Thi\`es. \ifkey\else\hfill\rule{2.2cm}{0.4pt}\fi
\item The number of hours worked, and the wages earned. \ifkey\else\hfill\rule{2.2cm}{0.4pt}\fi
\item The number of students absent in a class of 30, and the number of students present. \ifkey\else\hfill\rule{2.2cm}{0.4pt}\fi
\end{enumerate}
\begin{solution}
\textit{Scoring: 1 point per situation.}\\
(a) negative \quad (b) negative \quad (c) positive \quad (d) negative \quad (e) positive \quad (f) negative \quad (g) positive \quad (h) negative.
\end{solution}
\end{exo}

\Needspace{16\baselineskip}
\begin{center}\textsf{\small\bfseries\color{apnavy}Questions 5 to 8: find the average rate of change of each function over the interval $[-1,\,3]$.}\end{center}
\begin{multicols}{2}\raggedcolumns
\noindent\begin{minipage}{\linewidth}
\begin{exo}[Free response]{\pts{4}}
$f(x)=x^3-2x$
\ifkey\else\vspace{1.6cm}\fi
\begin{solution}
\textit{Scoring: 4 points: 1 for each function value, 2 for the average rate of change.}\\
$f(3)=21$, $f(-1)=1$: $\dfrac{21-1}{3-(-1)}=\dfrac{20}{4}=5$.
\end{solution}
\end{exo}
\end{minipage}\par\vspace{8pt}
\noindent\begin{minipage}{\linewidth}
\begin{exo}[Free response]{\pts{4}}
$g(x)=\dfrac{x+5}{x-4}$
\ifkey\else\vspace{1.6cm}\fi
\begin{solution}
\textit{Scoring: 4 points: 1 for each function value, 2 for the average rate of change.}\\
$g(3)=\frac{8}{-1}=-8$, $g(-1)=\frac{4}{-5}=-0.8$: $\dfrac{-8-(-0.8)}{4}=-1.8=-\dfrac95$.
\end{solution}
\end{exo}
\end{minipage}\par\vspace{8pt}
\noindent\begin{minipage}{\linewidth}
\begin{exo}[Free response]{\pts{4}}
\begin{center}\begin{tabular}{|c|c|}\hline$x$&$k(x)$\\\hline$-3$&$4$\\\hline$-1$&$7$\\\hline$1$&$-2$\\\hline$3$&$-5$\\\hline$5$&$0$\\\hline\end{tabular}\end{center}
\ifkey\else\vspace{0.6cm}\fi
\begin{solution}
\textit{Scoring: 4 points: 1 for each function value, 2 for the average rate of change.}\\
$k(3)=-5$ and $k(-1)=7$: $\dfrac{-5-7}{3-(-1)}=\dfrac{-12}{4}=-3$.
\end{solution}
\end{exo}
\end{minipage}\par\vspace{8pt}
\noindent\begin{minipage}{\linewidth}
\begin{exo}[Free response]{\pts{4}}
The graph of $q$:
\gr{xmin=-3.5,xmax=5.5,ymin=-2,ymax=4,width=6.2cm,height=4.6cm}{\addplot[apblue,very thick] coordinates {(-3,3) (1,-1) (3,-1)};\addplot[apblue,very thick] coordinates {(3,2) (5,0)};\pt{-3,3}\pt{5,0}\op{3,-1}\op{3,2}\pt{3,3}}
\ifkey\else\vspace{0.8cm}\fi
\begin{solution}
\textit{Scoring: 4 points: 1 for each function value, 2 for the average rate of change.}\\
$q(-1)=1$ and $q(3)=3$ (the filled point): $\dfrac{3-1}{3-(-1)}=\dfrac24=\dfrac12$.
\end{solution}
\end{exo}
\end{minipage}\par\vspace{8pt}
\noindent\begin{minipage}{\linewidth}
\begin{center}\textsf{\small\bfseries\color{apnavy}Questions 9 and 10: find the constant from the given average rate of change.}\end{center}
\begin{exo}[Free response]{\pts{5}}
The function $f$ is defined by $f(x)=kx^2-2x$, where $k$ is a constant. The average rate of change of $f$ over $[2,4]$ is $16$. Find $k$.
\ifkey\else\vspace{1.6cm}\fi
\begin{solution}
\textit{Scoring: 5 points: 2 for the function values, 2 for the equation, 1 for the value of $k$.}\\
$f(4)=16k-8$, $f(2)=4k-4$: $\dfrac{12k-4}{2}=6k-2=16$, so $k=3$.
\end{solution}
\end{exo}
\end{minipage}\par\vspace{8pt}
\noindent\begin{minipage}{\linewidth}
\begin{exo}[Free response]{\pts{5}}
\begin{center}\begin{tabular}{|c|c|}\hline$x$&$w(x)$\\\hline$1$&$6$\\\hline$5$&$m$\\\hline$7$&$3$\\\hline\end{tabular}\end{center}
The average rate of change of $w$ over $[1,5]$ is $3$. Find the constant $m$.
\ifkey\else\vspace{1.4cm}\fi
\begin{solution}
\textit{Scoring: 5 points: 3 for the equation, 2 for the value of $m$.}\\
$\dfrac{m-6}{5-1}=3$, so $m-6=12$ and $m=18$.
\end{solution}
\end{exo}
\end{minipage}\par\vspace{8pt}

\end{multicols}

\Needspace{16\baselineskip}
\begin{center}\textsf{\bfseries\color{apnavy}Part B \textbullet\ Multiple choice \textbullet\ No graphing calculator \textbullet\ 16 questions, 3 points each (48 points)}\end{center}
\begin{multicols}{2}\raggedcolumns
\noindent\begin{minipage}{\linewidth}
\gr{xmin=-6.5,xmax=6.5,ymin=-4,ymax=4,width=7.6cm,height=5cm}{\addplot[apblue,very thick,domain=-6:6,samples=150]{3*cos(deg(pi*x/6))};\node[font=\small,apblue,fill=white,inner sep=1pt] at (axis cs:-4.6,1.6){$y=g(x)$};}
\begin{center}\textsf{\small\bfseries\color{apnavy}Question 11 refers to the graph of $g$ on $-6\le x\le6$.}\end{center}
\begin{exo}{\pts{3}}
On which interval is $g$ increasing while its rate of change is decreasing?
\choices{$(-6,-3)$}{$(-3,0)$}{$(0,3)$}{$(3,6)$}
\answer{B}{On $(-3,0)$ the graph rises toward its maximum and is concave down, so the (positive) rate of change decreases.}
\end{exo}
\end{minipage}\par\vspace{8pt}
\noindent\begin{minipage}{\linewidth}
\begin{exo}{\pts{3}}
\par\noindent\begin{minipage}[t]{0.27\linewidth}\vspace{0pt}\centering\small\begin{tabular}{|c|c|}\hline$x$&$f(x)$\\\hline$0$&$4$\\\hline$2$&$10$\\\hline$3$&$13$\\\hline$5$&$19$\\\hline$9$&$31$\\\hline\end{tabular}\end{minipage}\hfill\begin{minipage}[t]{0.7\linewidth}\vspace{0pt}The table gives values of $f$. Which claim and explanation best fit these data?\end{minipage}\par
\choices{$f$ is best modeled by a linear function because the average rate of change over each interval is constant.}{$f$ is best modeled by a quadratic function because the average rate of change over each interval changes linearly.}{$f$ is best modeled by a quadratic function because the average rate of change over each interval is constant.}{$f$ is best modeled by a linear function because the average rate of change over each interval changes linearly.}
\answer{A}{The intervals have different lengths, but each average rate of change is $3$: $\frac62,\frac31,\frac62,\frac{12}4$. A constant rate of change means linear.}
\end{exo}
\end{minipage}\par\vspace{8pt}
\noindent\begin{minipage}{\linewidth}
\begin{exo}{\pts{3}}
\par\noindent\begin{minipage}[t]{0.27\linewidth}\vspace{0pt}\centering\small\begin{tabular}{|c|c|}\hline$x$&$k(x)$\\\hline$0$&$1$\\\hline$2$&$3$\\\hline$4$&$9$\\\hline$6$&$19$\\\hline$8$&$33$\\\hline\end{tabular}\end{minipage}\hfill\begin{minipage}[t]{0.7\linewidth}\vspace{0pt}The table gives values of $k$. Which claim and explanation best fit these data?\end{minipage}\par
\choices{$k$ is best modeled by a linear function because the rates of change over consecutive equal-length intervals are constant.}{$k$ is best modeled by a quadratic function because the rates of change over consecutive equal-length intervals are constant.}{$k$ is best modeled by a quadratic function because the rates of change over consecutive equal-length intervals change by a constant amount.}{$k$ is best modeled by a linear function because the rates of change over consecutive equal-length intervals change by a constant amount.}
\answer{C}{The rates over the intervals of length 2 are $1,3,5,7$: they increase by $2$ each time, which is the pattern of a quadratic function.}
\end{exo}
\end{minipage}\par\vspace{8pt}
\noindent\begin{minipage}{\linewidth}
\begin{exo}{\pts{3}}
\par\noindent\begin{minipage}[t]{0.27\linewidth}\vspace{0pt}\centering\small\begin{tabular}{|c|c|}\hline$x$&$m(x)$\\\hline$1$&$2$\\\hline$2$&$10$\\\hline$3$&$14$\\\hline$4$&$16$\\\hline$5$&$17$\\\hline\end{tabular}\end{minipage}\hfill\begin{minipage}[t]{0.7\linewidth}\vspace{0pt}The table gives values of $m$. Which claim and explanation best fit these data?\end{minipage}\par
\choices{The graph of $m$ could be concave up because the rates of change over consecutive equal-length intervals are positive.}{The graph of $m$ could be concave up because the rates of change over consecutive equal-length intervals are increasing.}{The graph of $m$ could be concave down because the rates of change over consecutive equal-length intervals are decreasing.}{The graph of $m$ could be concave down because the rates of change over consecutive equal-length intervals are positive.}
\answer{C}{The rates are $8,4,2,1$: decreasing, so concave down. Being positive tells us $m$ increases, not its concavity.}
\end{exo}
\end{minipage}\par\vspace{8pt}
\noindent\begin{minipage}{\linewidth}
\begin{exo}{\pts{3}}
\par\noindent\begin{minipage}[t]{0.27\linewidth}\vspace{0pt}\centering\small\begin{tabular}{|c|c|}\hline$x$&$n(x)$\\\hline$0$&$40$\\\hline$2$&$24$\\\hline$4$&$14$\\\hline$6$&$8$\\\hline$8$&$6$\\\hline\end{tabular}\end{minipage}\hfill\begin{minipage}[t]{0.7\linewidth}\vspace{0pt}The table gives values of $n$. Which claim and explanation best fit these data?\end{minipage}\par
\choices{The graph of $n$ could be concave up because the rates of change over consecutive equal-length intervals are increasing.}{The graph of $n$ could be concave up because the rates of change over consecutive equal-length intervals are negative.}{The graph of $n$ could be concave down because the rates of change over consecutive equal-length intervals are negative.}{The graph of $n$ could be concave down because the rates of change over consecutive equal-length intervals are decreasing.}
\answer{A}{The rates are $-8,-5,-3,-1$: increasing, so concave up. Negative rates only show that $n$ decreases.}
\end{exo}
\end{minipage}\par\vspace{8pt}
\noindent\begin{minipage}{\linewidth}
\begin{exo}{\pts{3}}
The figure below shows the graph of a function $f$. Which statement best describes $f$?
\sk{width=4.6cm,height=3.5cm,xmin=-0.3,xmax=5,ymin=-0.3,ymax=4.2}{\addplot[black,very thick,domain=0:3.8,samples=100]{4-0.25*x^2};}
\choices{The rate of change of $f$ is positive and increasing.}{The rate of change of $f$ is positive and decreasing.}{The rate of change of $f$ is negative and increasing.}{The rate of change of $f$ is negative and decreasing.}
\answer{D}{The graph falls (negative rate) and is concave down: it falls more and more steeply, so the rate decreases.}
\end{exo}
\end{minipage}\par\vspace{8pt}
\noindent\begin{minipage}{\linewidth}
\begin{exo}{\pts{3}}
The figure below shows the graph of a function $h$. Which statement about $h$ is correct?
\sk{width=4.6cm,height=3.5cm,xmin=-0.3,xmax=4.4,ymin=-0.3,ymax=6}{\addplot[black,very thick,domain=0.1:4,samples=100]{5*exp(-0.7*x)};}
\choices{$h$ is increasing, and its rate of change is increasing.}{$h$ is increasing, and its rate of change is decreasing.}{$h$ is decreasing, and its rate of change is increasing.}{$h$ is decreasing, and its rate of change is decreasing.}
\answer{C}{Falling and concave up: the negative rate of change increases toward $0$.}
\end{exo}
\end{minipage}\par\vspace{8pt}
\noindent\begin{minipage}{\linewidth}
\begin{exo}{\pts{3}}
The figure below shows the graph of a function $k$. Which statement about $k$ is correct?
\sk{width=4.6cm,height=3.5cm,xmin=-0.3,xmax=3,ymin=-0.3,ymax=4.4}{\addplot[black,very thick,domain=0.1:2.8,samples=100]{0.3+0.5*x^2};}
\choices{$k$ is increasing and the graph of $k$ is concave up.}{$k$ is increasing and the graph of $k$ is concave down.}{$k$ is decreasing and the graph of $k$ is concave up.}{$k$ is decreasing and the graph of $k$ is concave down.}
\answer{A}{The graph rises more and more steeply: increasing and concave up.}
\end{exo}
\end{minipage}\par\vspace{8pt}
\noindent\begin{minipage}{\linewidth}
\begin{exo}{\pts{3}}
The figure below shows the graph of a function $g$. Which statement is correct?
\sk{width=4.6cm,height=3.5cm,xmin=-0.3,xmax=4.6,ymin=-3.4,ymax=0.6}{\addplot[black,very thick,domain=0:4.3,samples=100]{-3+1.4*sqrt(x)};}
\choices{The rate of change of $g$ is positive and decreasing.}{The rate of change of $g$ is positive and increasing.}{The rate of change of $g$ is negative and increasing.}{The rate of change of $g$ is negative and decreasing.}
\answer{A}{The graph rises (positive rate) and flattens out, concave down (the rate decreases).}
\end{exo}
\end{minipage}\par\vspace{8pt}
\noindent\begin{minipage}{\linewidth}
\begin{center}\textsf{\small\bfseries\color{apnavy}Questions 20 to 23 refer to the graph of $g$. A and E are endpoints, B and D are relative extrema and C is the point of inflection; the letters stand for $x$-coordinates.}\end{center}
\sk{width=4.6cm,height=3.5cm,xmin=-0.3,xmax=6.6,ymin=0,ymax=16,width=6.6cm,height=4.4cm}{\addplot[black,very thick,domain=0:6,samples=100]{x^3/3-3*x^2+5*x+4+6};\lab{0}{10}{A}{east}\lab{1}{12.333}{B}{south}\lab{3}{7}{C}{south west}\lab{5}{1.667}{D}{north}\lab{6}{4}{E}{west}}
\begin{exo}{\pts{3}}
On which interval is $g$ decreasing and the graph of $g$ concave up?
\choices{From A to B}{From B to C}{From C to D}{From D to E}
\answer{C}{From the inflection point C to the minimum D, the graph falls and bends upward.}
\end{exo}
\end{minipage}\par\vspace{8pt}
\noindent\begin{minipage}{\linewidth}
\begin{exo}{\pts{3}}
On which interval is the rate of change of $g$ positive and decreasing?
\choices{From A to B}{From B to C}{From C to D}{From D to E}
\answer{A}{From A to the maximum B, $g$ increases and is concave down.}
\end{exo}
\end{minipage}\par\vspace{8pt}
\noindent\begin{minipage}{\linewidth}
\begin{exo}{\pts{3}}
On which interval is $g$ increasing and the graph of $g$ concave up?
\choices{From A to B}{From B to C}{From C to D}{From D to E}
\answer{D}{From the minimum D to the endpoint E.}
\end{exo}
\end{minipage}\par\vspace{8pt}
\noindent\begin{minipage}{\linewidth}
\begin{exo}{\pts{3}}
On which interval is the rate of change of $g$ negative and decreasing?
\choices{From A to B}{From B to C}{From C to D}{From D to E}
\answer{B}{From the maximum B to the inflection point C, $g$ decreases and is concave down.}
\end{exo}
\end{minipage}\par\vspace{8pt}
\noindent\begin{minipage}{\linewidth}
\begin{center}\textsf{\small\bfseries\color{apnavy}Questions 24 to 26 refer to the graph of $k$. A and E are endpoints, B and D are relative extrema and C is the point of inflection.}\end{center}
\gr{xmin=-5.8,xmax=5.8,ymin=-7,ymax=7,width=7.6cm,height=6.2cm}{\addplot[apblue,very thick,domain=-5:5,samples=150]{(x^3-27*x)/9};\pt{-5,1.111}\pt{-3,6}\pt{0,0}\pt{3,-6}\pt{5,-1.111}\node[font=\small\bfseries,anchor=north east,fill=white,inner sep=1pt] at (axis cs:-5,1.111){A};\node[font=\small\bfseries,anchor=south,fill=white,inner sep=1pt] at (axis cs:-3,6.2){B};\node[font=\small\bfseries,anchor=south west,fill=white,inner sep=1pt] at (axis cs:0.2,0.2){C};\node[font=\small\bfseries,anchor=north,fill=white,inner sep=1pt] at (axis cs:3,-6.2){D};\node[font=\small\bfseries,anchor=north west,fill=white,inner sep=1pt] at (axis cs:5,-1.111){E};}
\begin{exo}{\pts{3}}
Which claim and explanation is correct?
\choices{On the interval from B to C, $k$ is decreasing because the graph of $k$ is concave down.}{On the interval from B to C, $k$ is decreasing because the rate of change of $k$ is negative.}{On the interval from B to C, $k$ is decreasing because the rate of change of $k$ is decreasing.}{On the interval from B to C, $k$ is decreasing because the graph of $k$ is positive.}
\answer{B}{A function decreases where its rate of change is negative. Concavity describes how the rate changes, not its sign.}
\end{exo}
\end{minipage}\par\vspace{8pt}
\noindent\begin{minipage}{\linewidth}
\begin{exo}{\pts{3}}
Which claim and explanation is correct?
\choices{On the interval from D to E, the graph of $k$ is concave up because the rate of change of $k$ is positive.}{On the interval from D to E, the graph of $k$ is concave down because the rate of change of $k$ is increasing.}{On the interval from D to E, the graph of $k$ is concave up because the rate of change of $k$ is increasing.}{On the interval from D to E, the graph of $k$ is concave down because the rate of change of $k$ is positive.}
\answer{C}{Concave up means the rate of change is increasing.}
\end{exo}
\end{minipage}\par\vspace{8pt}
\noindent\begin{minipage}{\linewidth}
\begin{exo}{\pts{3}}
Which claim and explanation is correct?
\choices{On the interval from C to D, the rate of change of $k$ is decreasing because $k$ is negative.}{On the interval from C to D, the rate of change of $k$ is positive because the graph of $k$ is concave up.}{On the interval from C to D, the rate of change of $k$ is increasing because $k$ is decreasing.}{On the interval from C to D, the rate of change of $k$ is increasing because the graph of $k$ is concave up.}
\answer{D}{Between the inflection point C and the minimum D, the graph is concave up: the (negative) rate of change increases toward $0$.}
\end{exo}
\end{minipage}\par\vspace{8pt}

\end{multicols}
